%\documentclass[portrait,a0,draft]{a0poster}
\documentclass[portrait,a0]{a0poster}
%\usepackage{german,epsf,psfig,pstricks,multicol,pst-grad,color}
\usepackage{german,pstricks,multicol,pst-grad,color}
\usepackage{graphicx}

\begin{document}

% Fliessenden Hintergrund von RGB-Farbe 1. .98 .98 nach 1. .85 .85 
% und wieder nach  1. .98 .98 (1. .85 .85 wird nach 0.1=10% des Hinter-
% grunds angenommen)
% Achtung Werte unter .8 verbrauchen zu viel Tinte!!!
%\background{1. .98 .98}{1. .85 .85}{0.1}
% Hier haben wir ein paar Experimente gemacht und uns schliesslich auf
\background{.98 .98 1.}{0.698 .819 0.941}{0.05}
% geeinigt. Andere Versuche:
%\background{.98 .98 1.}{.0 .39 0.65}{0.1}
%\background{.98 .98 1.}{.0 .57 0.95}{0.1}
%\background{.98 .98 1.}{0.741 .870 0.941}{0.05}
%\background{.85 .85 1.}{.98 .98 1.}{-0.1}

% Groesse der einzelnen Spalten als Anteil der Gesamt-Textbreite
\renewcommand{\columnfrac}{.3}


% Posterueberschrift, optional mit RGB-Hintergrundfarbe (hier .99 .99 1.)
\begin{header}[1. .98 .98]
\begin{tabular}{|p{9.5in}|p{19.5in}|}
\hline
\begin{minipage}[t]{9.5in}
 \vspace*{0mm}
 \begin{center}
 \includegraphics[width=7in]{rwthaachenlogo.eps}
 \end{center}

 \vspace*{-0.8ex}
 \rule{9.5in}{0.05ex}

 \vspace*{0.6ex}
 \begin{tabular}{cp{5.5in}}
 \sf\veryHuge \raisebox{-1.05ex}{GRK 775} &
 \renewcommand{\baselinestretch}{0.8}\sf\Large
 Hierarchie und Symmetrie 
 in mathematischen Modellen
 \end{tabular}
 \renewcommand{\baselinestretch}{1}
 \vspace*{0.3ex}
\end{minipage}
&
\begin{minipage}[t]{19.5in}
 \vspace*{0mm}
\begin{center}
\vspace*{2ex}
\Huge Hier kommt der Titel rein

\vspace*{2ex}
\large Gerhard Hi\ss, Lehrstuhl D f\"ur Mathematik
\end{center}
\end{minipage} \\
\hline
\end{tabular}

\end{header}

%%% Beginn des Posters ...
\begin{poster}
%%% Beginn der ersten Poster-Spalte
\begin{pcolumn}
%%%%% Beginn einer Poster-Box, 
%%%%%   die Groesse wird an den eingegebenen Text angepasst!
      \begin{pbox}
        {\color{red} Pbox1 \vspace{35cm}}
      \end{pbox}
%%%%% Beginn der zweiten Poster-Box, diesmal mit RGB-Hintergrundfarbe
      \begin{pbox}[1. .98 .98]
        {\color{red} Pbox2 \vspace{37cm}}
      \end{pbox}
\end{pcolumn}
%%% Beginn der zweiten Poster-Spalte. 
%%%   Achtung, zw. \end{pccolumn} und \begin{pcolumn} keine Leerzeilen!   
\begin{pcolumn}
%%%%% Beginn einer Poster-Box, ohne Schatten
      \begin{pboxns}
       {\color{red} Pbox3 \vspace{35cm}}
      \end{pboxns}
      \begin{pbox}
        {\color{red} Pbox2 \vspace{35cm}}
      \end{pbox}
\end{pcolumn}
\begin{pcolumn}
\section{Ein Abschnitt}

Hier gibt es keine boxes.
\end{pcolumn}

\end{poster}

\end{document}

